\section{Introduction}

Cooperative multi-agent systems (MAS) are systems where multiple agents (such as autonomous vehicles/drones) work together to achieve a common goal such as maximizing a shared utility \citep{dorri2018multi}. These agents can communicate and coordinate with each other, either directly or indirectly, to solve complex tasks that are beyond a single agent's capabilities. Examples are drone swarms for autonomous aerial surveillance or disaster relief operations, or teams of robots that collaborate to explore unknown environments, harvest data from sensors, or manipulate objects, among others \citep{xie2017multi}.s.

Centralized approaches for addressing the MAS planning problem do not scale with the number of agents, as the amount of computational resources required to solve it quickly becomes prohibitive. In addition, the amount of information exchange required for a centralized planner to manage all agents can be unfeasibly high in large-scale settings, particularly in remote areas and disaster scenarios \citep{trodden2009robust}. Thus, recent research has focused on decentralized approaches for online MAS planning~\citep{zhang2021multi}. Indeed, they offer enhanced robustness, reduced computational burden, and lower communication load, particularly on infrastructure-based communications such as cellular radio access networks ~\citep{claes2017decentralised}.. 

The main challenge in decentralized approaches for cooperative MAS is optimizing agents' actions in a distributed manner to maximize a global reward function. This problem can typically be modeled as a Decentralized Partially Observable Markov Decision Process (Dec-POMDP)~\citep{oliehoek2016concise}. However, the computational complexity of Dec-POMDPs presents a significant hurdle, making direct optimal solution search infeasible in polynomial-time \citep{bernstein2002complexity}.
Several sampling-based planning algorithms have also been proposed to improve computational efficiency, particularly for special classes of Dec-POMDPs, such as multi-robot active perception \citep{zhang2021multi}.
A first family of approaches to address this is given by point-based methods, which scale well, but they may not cover the entire belief space well \citep{POINT_pineau2003point,POINT_shani2013survey}. Another set of algorithms is based on policy search \citep{POLICY_seuken2011memory, POLICY_amato2010optimizing} based on a parameterized policy representation. They can handle problems with large action spaces, but they may get stuck in local optima or require many samples
\citep{POLICY_amato2010optimizing}.
Thus, attention has turned towards algorithms based on Monte Carlo tree search (MCTS), due to their ability to effectively explore long planning horizons, their anytime nature~\citep{kocsis2006improved}, and their excellent performance in decentralized settings \citep{claes2017decentralised, best2019dec, czechowski2020decentralized}, effectively overcoming the limitations of other approaches.

In many present-day MAS application scenarios, the departure of agents from the system (henceforth denoted as \textit{attrition}, and due to e.g. failures or energy depletion) is a very common feature. However, all of the main approaches to decentralized MAS planning assume agents are always available and actively contributing to the joint planning process. When applied to scenarios with attrition, they perform in a heavily suboptimal manner and they often do not converge at all \citep{cybenko2021attritable,nguyen2022multi}. In swarm robotics, for instance, agent attrition due to robot failures, damage, or energy depletion affects the overall swarm behavior and task completion. Designing robust algorithms for decentralized MAS planning capable of effectively handling agent loss is critical for their successful deployment many in practical scenarios.

The common approaches for scenarios with attrition are based on periodically resetting agents' learned behavior, and restarting the learning process~\citep {avner2014concurrent,rosenski2016multi,hanawal2021multiplayer}. In volatile settings with frequent failures, such a feature may significantly hamper the overall performance of the active perception task, by slowing down the convergence rate and by keeping the system far from adapting and thus from achieving optimal operating conditions. 
Therefore, how to efficiently and effectively perform online MAS planning in the presence of attrition, while achieving fast convergence, is a key open issue.

In this paper, we develop a novel decentralized planning algorithm that achieves both of these objectives. Our approach is based on MCTS and the use of the global reward instead of the local one in the estimation of each agent's local contribution. Moreover, it exploits regret matching (RM)~\citep{hart2000simple} to coordinate the actions between agents. We prove that our approach guarantees that the average joint action of all agents converges to a Nash equilibrium (NE) if every agent applies the same RM procedure in any cooperative game with a submodular utility function. Arriving at an NE guarantees no diverging interest between the agents, and therefore it ensures that all participants come into a self-enforcing agreement to effectively coordinate their action decisions in a decentralized manner.
The main contributions of our work are:

\begin{itemize}[leftmargin=*]
\item We develop Attritable MCTS (A-MCTS), a new online decentralized planning algorithm, based on MCTS and Regret Matching, that can quickly and efficiently adapt the plan to settings where agents fail, even at high rates.
\item We show that, by modulating the utility function for each agent, under the assumption of submodularity, successive iterations are guaranteed to improve joint policies, and eventually lead to convergence of our algorithm.
We prove a strong convergence result for approximating a pure-strategy Nash equilibrium in a fully distributed fashion.
\item We evaluate our proposed approach in several information-gathering scenarios with attrition. Results suggest that, in the presence of frequent failures, our solution improves substantially over the best existing approaches in terms of global utility and scalability.
\end{itemize}